\documentclass[11pt,compress,t,notes=noshow, xcolor=table]{beamer}

\input{../../style/preamble}
\input{../../latex-math/basic-math}
\input{../../latex-math/basic-ml}

\title{Optimization in Machine Learning}

\comment[NOTE]{BB}{Beim soft-thresholding op: die connections / Äquivalenz zu Subgradienten hast du auf der Tonspur gesagt oder will man das den Leuten eher nicht zumuten?
Motivationsslide am Ende: "Warum verwendet in der Statistik -> weil in der Regel additive Modelle und damit kann man exakte Updates herleiten?"}

\begin{document}

\titlemeta{}{
Coordinate descent
}{
figure_man/Coordinate_descent.png
}{
\item Axes as descent direction
\item CD on linear model and Lasso
\item Soft thresholding
}

\begin{framei}{coordinate descent}
\item Assumption: objective function not differentiable
\item Idea: instead of gradient, use coordinate directions for descent
\item First: Select starting point $\bm{x}^{[0]} = (x^{[0]}_1, \ldots, x^{[0]}_d)$
\item Step $t$: Minimize $f$ along $x_i$ for each dimension $i$ for fixed $x^{[t]}_1, \ldots,x^{[t]}_{i-1}$ and $x^{[t-1]}_{i+1}, \ldots,x^{[t-1]}_d$:
\imageC[0.42][https://en.wikipedia.org/wiki/Coordinate_descent]{figure_man/Coordinate_descent.png}
\comment[NOTE]{BB}{austauschen mit eigener grafik}
\end{framei}


\begin{framev}{coordinate descent}
\begin{itemize}
\item Minimum is determined with (exact/inexact) line search
\item Order of dimensions can be any permutation of $\left\{1,2,\ldots,d\right\}$
\comment[NOTE]{CK}{verstehe ich nicht}
\item Convergence:
\comment[NOTE]{UU}{Diskussion von 3 Fällen nach Ryan Tibshirani.
1. f konvex + diff'bar
2. f konvex (Gegenbeispiel)
3. f Summe aus konvexer + diff'barer Funktion und separabler Funktion mit konvexen Summanden.
https://www.stat.cmu.edu/~ryantibs/convexopt-S15/lectures/22-coord-desc.pdf, Folien 5--7
Nachtrag: L(x) + ||x||_1 als ML-Beispiel hinzufügen.}
\begin{itemize}
\item $f$ convex differentiable
\comment[NOTE]{CK}{haben oben angenommen dass f non-differentiable}
\item $f$ sum of convex differentiable and \textit{convex separable} function:
\begin{equation*}
f(\xv) = g(\xv) + \sum_{i=1}^d h_i(x_i),
\end{equation*}
\item where $g$ convex differentiable and $h_i$ convex
\end{itemize}
\end{itemize}
\end{framev}


\begin{framev}{coordinate descent}
No convergence in general for convex functions, counterexample:
\imageC[0.6][https://en.wikipedia.org/wiki/Coordinate_descent]{figure_man/Nonsmooth_coordinate_descent.png}
\end{framev}


\begin{framev}{Example: linear regression}
Minimize LM with L2 loss via CD
$$\min g(\thetav) = \min_{\thetav} \frac{1}{2}\sumin \left(y^{(i)} - \thetav^\top \xi\right)^2 = \min_{\thetav} \frac{1}{2}\|\yv - \Xmat \thetav\|^2 $$
where $\yv \in \R^n$, $\Xmat \in \R^{n \times d}$ with columns $\xv_1, \ldots, \xv_d \in \R^n$\\
\spacer
Assume orthogonal design for intuition, i.e., $\Xmat^\top \Xmat = I_d$:
\begin{align*}
g(\thetav) &= \frac{1}{2}\yv^\top \yv + \frac{1}{2}\thetav^\top \thetav - \yv^\top \Xmat \thetav  \\
&\overset{(*)}{=} \frac{1}{2}\yv^\top \yv + \frac{1}{2}\thetav^\top \thetav - \yv^\top \sum_{k = 1}^d \xv_k \theta_k
\end{align*}
$^{(*)}$ $\Xmat \thetav = \xv_1 \theta_1 + \xv_2 \theta_2 + \cdots + \xv_d \theta_d = \sum_{k = 1}^d \xv_k \theta_k$
\end{framev}


\begin{framei}{example: linear regression}
\item Exact CD update in direction $j$:
$$\frac{\partial{g}(\thetav)}{\partial \theta_j} = \theta_j - \yv^\top \xv_j$$
\item By solving $\frac{\partial{g}(\thetav)}{\partial \theta_j} = 0$ we get $$\theta_j^\ast = \yv^\top \xv_j$$
\item Repeat this update for all $\theta_j$
\end{framei}


\begin{framev}{Soft thresholding}
Minimize LM with L2 loss and L1 regularization via CD
$$\min_{\thetav} h(\thetav) = \min_{\thetav} \frac{1}{2}\|\yv - \Xmat \thetav\|^2 + \lambda \|\thetav\|_1$$
Note that $h(\thetav) = \frac{1}{2}\yv^\top \yv + \frac{1}{2}\thetav^\top \thetav  - \sum_{k = 1}^d (\yv^\top \xv_k \theta_k + \lambda|\theta_k|)$\\
Again assume $\Xmat^\top \Xmat = I_d$. Since $|\cdot|$ not differentiable, distinguish cases:
\spacer
{\footnotesize
\begin{itemize}
\item \textbf{Case 1:} $\theta_j > 0$.
Then $|\theta_j| = \theta_j$ and
$$0 = \frac{\partial{h}(\thetav)}{\partial \theta_j} = \theta_j - \yv^\top \xv_j + \lambda \qquad \Leftrightarrow \qquad \theta^\ast_{j, \text{Lasso}} = \theta^\ast_j - \lambda$$
\item \textbf{Case 2:} $\theta_j < 0$.
Then $|\theta_j| = - \theta_j$ and
$$0 = \frac{\partial{h}(\thetav)}{\partial \theta_j} = \theta_j - \yv^\top \xv_j - \lambda \qquad \Leftrightarrow \qquad \theta^\ast_{j, \text{Lasso}} = \theta^\ast_j + \lambda$$
\item \textbf{Case 3:} $\theta_j = 0$
\end{itemize}}
\comment[NOTE]{BB}{ist die analyse hier wirklich völlig richtig? theta_j ist doch variabel und wir können nicht einfach...
naja es stimmt schon, hängt davon ab wo unser optimum liegt. vielleicht sollte man das mal zeichnen
also die analyse stimmt, zeichnen sollten wir das trotzdem, einmal wo das optimum im diffbaren ast l...
UU:
Ich befürchte, dass ein Bild auch nicht viel hilft. Man braucht Subgradienten zur Herleitung des Soft-Thresholding-Operators. Und Subgradienten werden nicht behandelt.}
\end{framev}


\begin{framev}{soft thresholding}
We can write the solution as:
$$\theta^\ast_{j, \text{Lasso}} =
\begin{cases}
\theta^\ast_j - \lambda & \text{ if } \theta^\ast_j > \lambda \\
\theta^\ast_j + \lambda & \text{ if } \theta^\ast_j < - \lambda \\
0 &  \text{ if }  \theta^\ast_j \in [- \lambda, \lambda],
\end{cases}$$
\comment[NOTE]{BB}{
man kann das in zusammenhang mit dem subgradienten bringen. der sieht ja genauso aus
hier ist übrigens eine übung von david dazu, die wie ansehen sollten
https://github.com/slds-lmu/lecture_cim2/blob/master/2020/Aufgaben/lasso.tex
}
This operation is called soft thresholding\\
\spacer
Coefs for which solution to unregularized problem is smaller than threshold ($|\thetav^\ast_j| < \lambda$) are shrunk to zero
\spacer
NB: Derivation of soft thresholding operator not trivial (subgradients)
\end{framev}


\begin{framev}{CD for statistics and ML}
Why is it being used?
\begin{itemizeL}
\item Easy to implement
\item Scalable: no storage/operations on large objects,
just current point \\
\comment[NOTE]{BB}{es wäre schön die effizienz von CD gegen GD oder anderes zu zeigen. das könnte auch in einem extra chucnk mal vergleichen.
so large scale betrachtungen}
$\Rightarrow$ Good implementation can achieve state-of-the-art performance
\item Applicable for non-differentiable (but convex separable) objectives
\end{itemizeL}
\spacer
Examples:
\begin{itemizeL}
\item Lasso regression, Lasso GLM, graphical Lasso
\comment[NOTE]{BB}{könnten hier am ende zb einen slide zu liblinear mal einabeun mit etwas info zu dessen runtiume auf einem grossen porblem}
\item Support Vector Machines
\item Regression with non-convex penalties
\end{itemizeL}
\end{framev}
\comment[NOTE]{BB}{weitere motivation: coord desc bei additiven modellen besprechen. besonders zusammenhang zu CWB}
\comment[NOTE]{BB}{wir sollten das hier auch besser visaluwerien für das LM und den lasso.
also den optim trace in 2d mit uniavr funs}

\endlecture
\end{document}
